\section{The three subtopoi of $\mathbb N$-sets}\label{sec:dynamics}

An action of the additive monoid $(\mathbb N,+)$ on a set $X$ is the same as a map $\tau : X\to X$ ($x\cdot n = \tau^n(x)$): a discrete dynamical system. A \emph{subtopos} of the topos of such systems is a full subcategory closed under limits, whose inclusion has a left adjoint preserving finite limits.

\begin{theorem}\label{th:dynamics}
The topos of sets with an endomorphism has exactly three subtopoi:
\begin{enumerate}
\item all of it;
\item the sets with an \emph{automorphism} ($\tau$ bijective, i.e.\ $\mathbb Z$-sets), the left adjoint being $X\mapsto\varinjlim(X\xrightarrow{\tau}X\xrightarrow{\tau}\cdots)$;
\item the degenerate subtopos (only the one-point set).
\end{enumerate}
Equivalently, the one-object category $\B\mathbb N$ carries exactly three Grothendieck topologies.
\end{theorem}

\begin{trace}
\item[deck] sets with an endomorphism $=$ presheaves on $\B\mathbb N$.
\item[arch] Bridge~\ref{br:sub}: subtopoi $\leftrightarrow$ Grothendieck topologies on $\B\mathbb N$.
\item[specialise] sieves on the object are $\emptyset$ and the right ideals $I_n = \{n,n+1,\dots\}$; pulling $I_n$ back along $m$ gives $I_{\max(n-m,0)}$. Count the topologies: three.
\item[unfold] identify the sheaves of each topology.
\item[verify] direct proof below; Lean: %LEAN-NAT%
being formalised.
\item[novelty] not found in print.
\end{trace}

\direct{} Let $J$ be a topology; it contains $I_0$. If it contains $\emptyset$, it contains every sieve (covering sieves are closed upwards). Otherwise, if it contains some $I_n$ with $n\ge 1$, it contains $I_1$ (the pullback of $I_n$ along $n-1$), and then every $I_k$ by induction: the pullback of $I_{k+1}$ along any $r\ge 1$ in the covering sieve $I_1$ is $I_{k+1-r}\supseteq I_k$, covering, so $I_{k+1}$ covers by transitivity. Hence three topologies: $\{I_0\}$, $\{I_n : n\ge 0\}$, all sieves. For the middle one, a matching family on $I_1$ is determined by its value $x_1$ at $1$ (then $x_{1+k} = \tau^k(x_1)$), and the sheaf condition asks for a unique $x$ with $\tau(x) = x_1$: $\tau$ is bijective. The left adjoint inverts $\tau$ by a filtered colimit, which preserves finite limits.\qed

\inwords{} There is no intermediate way of \q{localising} dynamical systems: one either keeps all of them, or makes the dynamics invertible by passing to its colimit, or collapses everything. Sets with a fixed point, eventually periodic systems or finite systems do not form subtopoi.

\begin{example}
For a group $G$, the only sieves on $\B G$ are $\emptyset$ and $G$, so $G$-sets have exactly two subtopoi. For $(\mathbb N^2,+)$ there are more than three: inverting the first generator only gives a fourth.
\end{example}
